% Beamer's default font size is 11 points. It is possible to set the
% default font size to any of 8, 9, 10, 11, 12, 14, 17, 20
\documentclass[13pt]{beamer}

\usetheme{Warsaw}
%\usecolortheme{dolphin}
\usecolortheme{seahorse}
%\useinnertheme{circles}
\usefonttheme{professionalfonts}
%\usefonttheme{structurebold}
%\usefonttheme{default}

\setbeamertemplate{navigation symbols}{}
 
\usepackage[english]{babel}
\usepackage[latin1]{inputenc}
\usepackage{times}
\usepackage[T1]{fontenc}
\usepackage{fancybox}

% These things can change:
\newcommand{\baseurl}{http://cell.vtt.fi/latex}
\urlstyle{sf}
\newcommand{\plus}{\makebox[2em]{+}}
\newcommand{\minus}{\makebox[2em]{-}}
\newcommand{\gray}[1]{%
\textcolor{gray}{#1}
}

% In beamer, defining new commands inside the slides fails.
% Definitions here:

\begin{document}

\logo{\includegraphics[width=1.4cm]{img/vttplain}\quad}
\title{\scalebox{1.5}{\textrm{\LaTeX}} Course 2011}
\subtitle{Part 4: Typesetting mathematics}
\author{Arho Virkki}
\institute{\textsc{VTT Technical Research Centre of Finland}}
\date{\scalebox{0.8}{\textcolor{gray}{}}}

\begin{frame}
  \titlepage
\end{frame}


\begin{frame}[fragile]\frametitle{Typesetting mathematics}
Inlined equations are written inside \verb|$$| characters:\\
E.g., \verb|$\sum_{k=1}^\infty a_k$| produces $\sum_{k=1}^\infty a_k$.\medskip

Equations on their own line are put into a separate
environment
\verb|\begin{equation}..\end{equation}|\footnote{This produces a number.},
\begin{equation}
\sum_{k=1}^\infty a_k,
\end{equation}
or \verb|\[..\]|\footnote{This doesn't.},
\[ \sum_{k=1}^\infty a_k. \] 
\end{frame}


\begin{frame}[fragile]\frametitle{Typesetting mathematics\dots}
Observe the differences in style when writing the equations within the text, or
separately (e.g. the limits of the sum were written differently)

\bigskip
The style can also be forced with the commands,
\begin{itemize}
\item \verb|\textstyle|
\item \verb|\displaystyle|
\end{itemize}
if needed.
\end{frame}


\begin{frame}[fragile]\frametitle{Typesetting mathematics\dots}
In the math mode, letters (the variables or constants) are typeset with slanted
typeface and extra spaces are ignored:
\begin{verbatim}
\begin{equation}
I(x) := a_1^2       b_{1,2}   c     x    .
\end{equation}
\end{verbatim}
\begin{equation}
I(x) := a_1^2       b_{1,2}   c     x    .
\end{equation}
\begin{flushleft}
If the equations are part ot the sentence,  commas should be written such that
the formulas become a natural part of the sentence.
\end{flushleft}
\end{frame}


\begin{frame}[fragile]\frametitle{Typesetting mathematics\dots}
The characters 
\begin{verbatim}
# $ % & ~ _ ^ \ { }
\end{verbatim}
have a special meaning in \LaTeX{}. If they are needed in the equations, 
one must write
\begin{verbatim}
\#   \$   \%   \&   \_   \{   \}
\end{verbatim}
For example,
\begin{verbatim}
\[ \mathcal{L}\{f\}(s) := 
   \int_0^\infty f(t) e^{-st} dt \]
\end{verbatim}
\[ \mathcal{L}\{f\}(s) := 
   \int_0^\infty f(t) e^{-st} dt \]
\end{frame}


\begin{frame}[fragile]\frametitle{Typesetting mathematics\dots}
\begin{itemize}
\item Because ordinary keyboards are lacking most of the mathematical symbols,
there are quite a number of special commands -- in addition to the greek
letters  $\alpha, \beta, \gamma, \delta,�\epsilon, \zeta, \eta, \theta,
\iota,�\kappa, \lambda, \mu, \nu, \xi, o, \pi, \rho, \sigma,
\tau, \upsilon, \phi, \chi, \psi$ and $\omega$\medskip
\item The commands for mathematical typesetting are best learned by heart.
Editors (including Eclipse) also help in remembering the commands.
\end{itemize}
\end{frame}


\begin{frame}[fragile]\frametitle{Upper and lower indices}
The upper and lower indices can be produced with \verb|^| and \verb|_|. \bigskip

\begin{verbatim}
\[ V^2_1 \ge V_{21} \]
\end{verbatim}

\[ V^2_1 \ge V_{21} \]
If there are more than one index, they need to be grouped with curly braces.
\end{frame}


\begin{frame}[fragile]\frametitle{Rational expressions}
\begin{verbatim}
\[ \frac{a}{b} \]
\[ \sqrt{2} \approx 1.4142 \]
\end{verbatim}\bigskip

\[ \frac{a}{b} \]
\[ \sqrt{2} \approx 1.4142 \]
\end{frame}


\begin{frame}[fragile]\frametitle{Rational expressions\dots}
\begin{verbatim}
\[ \left(
   \frac{\sqrt[3]{z-1}}{2}
   \right)^2 \]
\end{verbatim}\bigskip
\[ \left(\frac{\sqrt[3]{z-1}}{2}\right)^2 \]
\end{frame}


\begin{frame}[fragile]\frametitle{Rational expressions\dots}
Revision:
\begin{itemize}
\item Fraction: \verb|\frac{numerator}{denominator}| 
\item Root: \verb|\sqrt{}| ja \verb|\sqrt[n]{}| (!)
\item The size of the parenthesis can be automatically adjusted with the command
pair \verb|\left(| and \verb|\right)| 
\begin{itemize}
\item Also \verb|\left[, \left\{| and \verb+\left|+ work
\item One can mix $|, (, [, \{$ and $.$
\item The dot shows nothing. Example: $\displaystyle\left\{
\frac{\frac{a}{b}}{C}\right.$
\end{itemize}
\end{itemize}
\end{frame}


\begin{frame}[fragile]\frametitle{Rational expressions\dots}
Chain fractions (or continuous fractions):\bigskip

\begin{verbatim}
\[ a_0 + \cfrac{1}
         {a_1 + \cfrac{1}
                {a_2+\frac{1}{a_3}}} \]
\end{verbatim}\bigskip
\[ a_0 + \cfrac{1}
         {a_1 + \cfrac{1}
                {a_2+\frac{1}{a_3}}} \]
\end{frame}


\begin{frame}[fragile]\frametitle{Differential operators}
\begin{verbatim}
\[ \frac{dy}{dt} = f(y;\lambda) \]
\[ \frac{\partial u}{\partial t} = \kappa^2 
   \left( \frac{\partial^2 u}{\partial x^2}
        + \frac{\partial^2 u}{\partial y^2}
        + \frac{\partial^2 u}{\partial z^2} 
      \right) \]
\end{verbatim}
\[ \frac{dy}{dt} = f(y;\lambda) \]
\[ \frac{\partial u}{\partial t} = \kappa^2 
   \left( \frac{\partial^2 u}{\partial x^2}
        + \frac{\partial^2 u}{\partial y^2}
        + \frac{\partial^2 u}{\partial z^2} 
      \right) \]
\end{frame}


\begin{frame}[fragile]\frametitle{Limits, sums, integrals}
\begin{verbatim}
\[ \lim_{x \to 0} 
   \frac{3x^2 +7x^3}{x^2 +5x^4} = 3.\] 
\end{verbatim}
\[ \lim_{x \to 0} 
   \frac{3x^2 +7x^3}{x^2 +5x^4} = 3.\]\bigskip

\begin{verbatim}
\[ \sum_{k=1}^\infty \frac{1}{2^k} = 
   \frac{\pi^2}{6} \]
\end{verbatim}
\[ \sum_{k=1}^\infty \frac{1}{2^k} = 
   \frac{\pi^2}{6} \]
\end{frame}


\begin{frame}[fragile]\frametitle{Limits, sums, integrals\dots}
\begin{verbatim}
\[ \int_{-\infty}^\infty \sin(x)\ dx = \pi \]
\end{verbatim}
\[ \int_{-\infty}^\infty \sin(x)\ dx = \pi \]\bigskip
\begin{verbatim}
\[ \oint_{\partial C} h(z)\ dz = 0 \]
\end{verbatim}
\[ \oint_{\partial C} h(z)\ dz = 0 \]
\end{frame}


\begin{frame}[fragile]\frametitle{Limits, sums, integrals\dots}
\begin{verbatim}
\[ \int\limits_{-\infty}^\infty 
   \sin(x)\ dx = \pi \]
\end{verbatim}
\[ \int\limits_{-\infty}^\infty 
   \sin(x)\ dx = \pi \]\bigskip
\begin{verbatim}
\[ \oint\limits_{\partial C} h(z)\ dz = 0 \]
\end{verbatim}
\[ \oint\limits_{\partial C} h(z)\ dz = 0 \]
\end{frame}


\begin{frame}[fragile]\frametitle{Revision and some trivia}
\begin{itemize}
\item The upper and lower indices are obtained with \verb|^| ja \verb|_|
\item Spaces in equations can be enforced with \verb|\ |\\
(backslash and space)
\item The command \verb|\limits| puts more emphasis on the upper and lower
indices
\end{itemize}
\end{frame}


\begin{frame}[fragile]\frametitle{Spaces in equations}
\begin{verbatim}
\[ a\, b \]
\[ a\: b \]
\[ a\  b \]
\[ a\! b \]
\end{verbatim}
\[ a\, b \]
\[ a\: b \]
\[ a\  b \]
\[ a\! b \]
\begin{itemize}
\item The last command \verb|\!| is a negative space.
\item Use these fine-tunings sparingly!
\end{itemize}
\end{frame}


\begin{frame}[fragile]\frametitle{Math fonts}
Math mode is not for writing text. Nevertheless, there are some commands to
change the appearance of the fonts.\medskip

\begin{tabular}{ll}
\verb|A| & $A$ \\
\verb|\mathrm{A}| & $\mathrm{A}$ \\
\verb|\mathbf{A}| & $\mathbf{A}$ \\
\verb|\mathbb{A}| & $\mathbb{A}$  \\
\verb|\mathcal{A}| & $\mathcal{A}$, \\
\end{tabular}\medskip

The last two commands only work for capital letters.
\end{frame}


\begin{frame}[fragile]\frametitle{Font style, examples}
\begin{verbatim}
\[ \dot{M}_{\mathrm{CO_2}} \]
\end{verbatim}
\[ \dot{M}_{\mathrm{CO_2}} \]
\begin{verbatim}
\[ \alpha \in \mathbb{C}, \mathbf{v} \in X 
     \Rightarrow \alpha \mathbf{v} \in X \]
\end{verbatim}
\[ \alpha \in \mathbb{C}, \mathbf{v} \in X 
     \Rightarrow \alpha \mathbf{v} \in X \]
\begin{itemize}
\item Chemical formulas are written with capital letter.
\item (In modern literature, vectors are typically written just as $v \in X$,
without any additional decorations.
\end{itemize}
\end{frame}


\begin{frame}[fragile]\frametitle{Font size}
The size of the font is also controlled differently:\bigskip

\begin{tabular}{ll}
\verb|{\scriptscriptstyle \sum}| & ${\scriptscriptstyle \sum}$ \\
\verb|{\scriptstyle \sum}| & ${\scriptstyle \sum}$ \\
\verb|{\textstyle \sum}| & ${\textstyle \sum}$  \\
\verb|{\displaystyle \sum}| & ${\displaystyle \sum}$ \\
\end{tabular}
\end{frame}


\begin{frame}[fragile]\frametitle{Function names}
\LaTeX knows the most conventional function names:\bigskip
\begin{verbatim}
\[ \cos (\pi) = -1 \]
\end{verbatim}
\[ \cos (\pi) = -1 \]\bigskip

\begin{verbatim}
\[ \sin (0) = 0 \]
\end{verbatim}
\[ \sin (0) = 0 \]
\end{frame}


\begin{frame}[fragile]\frametitle{Function names\dots}
Known names include
\begin{verbatim}
\arccos \cos  \csc    \exp  \ker    \limsup 
\min    \sinh \arcsin \cosh \def    \gcd  
\lg     \ln   \Pr     \sup  \arctan \cot 
\det    \hom  \lim    \log  \sec    \tan
\arg    \coth \dim    \inf  \liminf \max    
\sin    \tanh
\end{verbatim}
\end{frame}


\begin{frame}[fragile]\frametitle{Function names\dots}
Own functions can be \textbf{declared} as operators with the command
\verb|\operatorname|.\bigskip
\begin{verbatim}
\[ \operatorname{arg min}_
   \Theta f(\Theta) \]
\end{verbatim}
\[ \operatorname{arg min}_\Theta f(\Theta) \]
\end{frame}

\begin{frame}[fragile]\frametitle{Function names\dots}
The previous example still needs some tuning:
\begin{verbatim}
\[ \operatorname{arg\ min}_
   \Theta f(\Theta) \]
\end{verbatim}
\[ \operatorname{arg\ min}_\Theta f(\Theta) \]
\begin{verbatim}
\[ \underset{\Theta}
   {\operatorname{arg\ min}} f(\Theta) \]
\end{verbatim}
\[ \underset{\Theta}{\operatorname{arg\ min}} f(\Theta) \]
\end{frame}


\begin{frame}[fragile]\frametitle{Mixing text and formulas}
LaTeX :
\begin{verbatim}
\[ S := \{ x\in \Omega | f(x) = c 
           \ \mathrm{and}\ g(x) < 0 \} \]
\end{verbatim}\bigskip
\[ S := \{ x\in \Omega | f(x) = c, \ \mathrm{and}\ g(x) < 0 \} \]
AMS-LaTeX:
\begin{verbatim}
\[ S := \{ x\in \Omega | f(x) = c
           \text{ and } g(x) < 0 \} \]
\end{verbatim}\bigskip
\[ S := \{ x\in \Omega | f(x) = c, \text{ and } g(x) < 0 \} \]
\end{frame}


\begin{frame}[fragile]\frametitle{Three dots(\dots)}
Three dots:\bigskip
\begin{verbatim}
\[ a_0 + a_2 + \cdots + a_n \]
\[ \ldots, \cdots, \ddots, \vdots \]
\end{verbatim} \bigskip
\[ a_0 + a_2 + \cdots + a_n \]
\[ \ldots, \cdots, \ddots, \vdots \]
\end{frame}


\begin{frame}[fragile]\frametitle{Mathematical accents}
\begin{tabular}{llll}
\verb|\underline{a}| & $\underline{a}$ &
\verb|\overline{a}| & $\overline{a}$ \\
\verb|\hat{a}| & $\hat{a}$ &
\verb|\check{a}| & $\check{a}$ \\
\verb|\tilde{a}| & $\tilde{a}$ &
\verb|\acute{a}| & $\acute{a}$ \\
\verb|\grave{a}| & $\grave{a}$ &
\verb|\dot{a}| & $\dot{a}$ \\
\verb|\ddot{a}| & $\ddot{a}$ &
\verb|\breve{a}| & $\breve{a}$ \\
\verb|\bar{a}| & $\bar{a}$ &
\verb|\vec{a}| & $\vec{a}$ 
\end{tabular}
\end{frame}


\begin{frame}[fragile]\frametitle{Norm and absolute value}
Norm:\bigskip
\begin{verbatim}
\[ \|a x\| = |a| \|x\| \]
\end{verbatim} \bigskip
\[ \|a x\| = |a| \|x\| \]
\end{frame}


\begin{frame}[fragile]\frametitle{Parentheses}
Dot produces an invisible parenthesis:\bigskip
\begin{verbatim}
\[ \left( \frac{1+z}{1-z} \right. \]
\[ \left. \frac{du}{dx} \right|_{x=0} \]
\end{verbatim}\bigskip
\[ \left( \frac{1+z}{1-z} \right. \]
\[ \left. \frac{du}{dx} \right|_{x=0} \]
\end{frame}


\begin{frame}[fragile]\frametitle{Multiline equations}
\begin{verbatim}
\begin{eqnarray}
F(x) &= & \int_a^b I(u,x) du  \nonumber \\
     &= & \frac{1}{\sqrt{\pi}}G(x)
\end{eqnarray}
\end{verbatim}\bigskip
\begin{eqnarray}
F(x) &= & \int_a^b I(u,x) du  \nonumber \\
     &= & \frac{1}{\sqrt{\pi}}G(x)
\end{eqnarray}
\end{frame}


\begin{frame}[fragile]\frametitle{Multiline equations\dots}
\begin{itemize}
\item \verb|eqnarray| works like a table, but it has three columns fixed.
\item The numbering can be suppressed with \verb|\nonumber|
\item If the numbering is not needed at all, one can use the form
\verb|\begin{eqnarray*}| \dots \verb|\end{eqnarray*}|
\end{itemize}
\end{frame}


\begin{frame}[fragile]\frametitle{Matrices}
Like the table,\bigskip
\begin{verbatim}
\[ \left[\begin{array}{ccc}
    1 & 0 & 2  \\
    0 & 2 & 0  \\
    0 & 0 & a \end{array}\right] \]
\end{verbatim}

\[ \left[\begin{array}{ccc}
    1 & 0 & 2  \\
    0 & 2 & 0  \\
    0 & 0 & a \end{array}\right] \]
\end{frame}


\begin{frame}[fragile]\frametitle{Matrices\dots}
Using the AMS-package:\bigskip
\begin{verbatim}
\[ \begin{bmatrix}
  1 & 0 & 2  \\
  0 & 2 & 0  \\
  0 & 0 & a \end{bmatrix} \]
\end{verbatim}
\[ \begin{bmatrix}
  1 & 0 & 2  \\
  0 & 2 & 0  \\
  0 & 0 & a \end{bmatrix} \]
\end{frame}


\begin{frame}[fragile]\frametitle{Matrices\dots}
\dots or \dots \bigskip
\begin{verbatim}
\[ \begin{pmatrix}
  1 & 0 & 2  \\
  0 & 2 & 0  \\
  0 & 0 & a \end{pmatrix} \]
\end{verbatim}
\[ \begin{pmatrix}
  1 & 0 & 2  \\
  0 & 2 & 0  \\
  0 & 0 & a \end{pmatrix} \]
\end{frame}


\begin{frame}[fragile]\frametitle{Functions defined in parts}
Like with tables,\bigskip
\begin{verbatim}
\[ |x| = \left\{ \begin{array}{rl}
    x & \mbox{if $x \geq 0$};\\
   -x & \mbox{if $x < 0$}.\end{array}
   \right. \] 
\end{verbatim}
\[ |x| = \left\{ \begin{array}{rl}
         x & \mbox{if $x \geq 0$};\\
        -x & \mbox{if $x < 0$}.\end{array} \right. \] 
\end{frame}


\begin{frame}[fragile]\frametitle{Functions defined in parts\dots}
Or using the AMS-package:\bigskip
\begin{verbatim}
\[ |x| = \left\{ \begin{aligned}
    x & \text{ if } x \geq 0;\\
   -x & \text{ if } x < 0.\end{aligned}
   \right. \] 
\end{verbatim}
\[ |x| = \left\{ \begin{aligned}
    x & \text{ if } x \geq 0;\\
   -x & \text{ if } x < 0.\end{aligned}
   \right. \] 
\end{frame}


\begin{frame}[fragile]\frametitle{Grouping with curly braces}
Equations can be annotated to help the reader:\bigskip
\begin{verbatim}
\[ \underbrace{A}_{\star} \]
\end{verbatim}
\[ \underbrace{A}_{\star} \]\bigskip
\begin{verbatim}
\[ \overbrace{B}^{\dag} \]
\end{verbatim}
\[ \overbrace{B}^{\dag} \]
\end{frame}



\begin{frame}[fragile]\frametitle{Grouping with curly braces\dots}
Example:\bigskip
\begin{verbatim}
\[ y' = \overbrace{
   \underbrace{Ay}_{\text{linear part}} + 
   \underbrace{By^R}_{\text{nonlinear part}}
   }^{\text{the driving force}} \]
\end{verbatim}
\[ y' = \overbrace{
        \underbrace{Ay}_{\text{linear part}} + 
        \underbrace{By^R}_{\text{nonlinear part}}
      }^{\text{the driving force}} \]
\end{frame}


\begin{frame}[fragile]\frametitle{Boxes}
Adding a box around an equation\bigskip
\begin{verbatim}
\[ \boxed{\frac{1}{1+x}} \]
\end{verbatim}
\[ \boxed{\frac{1}{1+x}} \]\bigskip

\begin{itemize}
\item may simplify the presentation
\item \dots or make it even more complicated
\end{itemize}
\end{frame}


\begin{frame}[fragile]\frametitle{Boxes\dots}
Example:\bigskip
\begin{verbatim}
\begin{eqnarray*}  
\| x - z \| &= &\| x\ \boxed{-y+y} - z \|\\
         &\leq &\| x - y \| + \| y - z \| 
\end{eqnarray*}
\end{verbatim}
\begin{eqnarray*}  
\| x - z \| &=    &\| x\ \boxed{-y+y} - z \| \\
            &\leq &\| x - y \| + \| y - z \| 
\end{eqnarray*}
\end{frame}


\begin{frame}[fragile]\frametitle{The AMS package}
In this course, we use the commands from the AMS package without separate
notification (as they have become standard over the years). The following
commands are especially from the AMS package \texttt{amsmath}:\bigskip

\begin{tabular}{ll}
\verb|\iint| & $\iint$ \\
\verb|\iiint| & $\iiint$ \\
\verb|\iiiint| & $\iiiint$ \\
\verb|\idotsint| & $\idotsint$ \\
\end{tabular}
\end{frame}


\begin{frame}[fragile]\frametitle{The AMS package\dots}
\begin{verbatim}
\[ \sum_{\substack{k=1\dots n\\ 
                   l=1\dots m\\
                   k\neq l}} 
    a_{k,l} \]
\end{verbatim}
\[ \sum_{\substack{k=1\dots n\\ 
                   l=1\dots m\\
                   k\neq l}} 
    a_{k,l} \]
\end{frame}


\begin{frame}[fragile]\frametitle{The AMS package\dots}
\begin{verbatim}
\[ \overset{*}{A}, \underset{*}{B} \]
\end{verbatim}
\[ \overset{*}{A}, \underset{*}{B} \]\bigskip
\begin{verbatim}
\[ \sideset{_a^b}{_c^d}\prod \]
\end{verbatim}
\[ \sideset{_a^b}{_c^d}\prod \]
\end{frame}


\begin{frame}[fragile]\frametitle{Greek letters}
\begin{tabular}{llllll}
\verb|\alpha| & $\alpha$ &
\verb|\beta| & $\beta$ &
\verb|\gamma| & $\gamma$ \\
\verb|\delta| & $\delta$  &
\verb|\epsilon| & $\epsilon$ &
\verb|\varepsilon| & $\varepsilon$\\
\verb|\zeta| & $\zeta$ &
\verb|\eta| & $\eta$ &
\verb|\theta| & $\theta$ \\
\verb|\vartheta| & $\vartheta$ &
\verb|\iota| & $\iota$ & 
\verb|\kappa| & $\kappa$ \\
\verb|\lambda| & $\lambda$ &
\verb|\mu| & $\mu$ &
\verb|\nu| & $\nu$ \\
\verb|\xi| & $\xi$ &
\verb|\pi| & $\pi$ &
\verb|\varpi| & $\varpi$\\
\verb|\rho| & $\rho$ &
\verb|\varrho| & $\varrho$ &
\verb|\sigma| & $\sigma$ \\
\verb|\varsigma| & $\varsigma$ &
\verb|\tau| & $\tau$ &
\verb|\upsilon| & $\upsilon$ \\
\verb|\phi| & $\phi$ &
\verb|\varphi| & $\varphi$ &
\verb|\chi| & $\chi$ \\
\verb|\psi| & $\psi$ &
\verb|\omega| & $\omega$ 
\end{tabular}
\end{frame}


\begin{frame}[fragile]\frametitle{Greek capital letters}
Only those letters are having a separate command, which do not have a
key in the ordinary keyboard.\bigskip

\begin{tabular}{llllll}
\verb|\Gamma| & $\Gamma$ &
\verb|\Delta| & $\Delta$ &
\verb|\Theta| & $\Theta$ \\
\verb|\Lambda| & $\Lambda$ &
\verb|\Xi| & $\Xi$ &
\verb|\Pi| & $\Pi$ \\
\verb|\Sigma| & $\Sigma$ &
\verb|\Upsilon| & $\Upsilon$ &
\verb|\Phi| & $\Phi$ \\
\verb|\Psi| & $\Psi$ &
\verb|\Omega| & $\Omega$
\end{tabular}
\end{frame}


\begin{frame}[fragile]\frametitle{Symbols: relations}
\begin{tabular}{llllll}
\verb|\approx| & $\approx$ &
\verb|\asymp| & $\asymp$ &
\verb|\bowtie| & $\bowtie$ \\
\verb|\cong| & $\cong$ &
\verb|\dashv| & $\dashv$ &
\verb|\doteq| & $\doteq$ \\
\verb|\equiv| & $\equiv$ &
\verb|\frown| & $\frown$ &
\verb|\ge| or \verb|\geq| & $\geq$ \\
\verb|\gg| & $\gg$ &
\verb|\in| & $\in$ &
\verb|\le| or \verb|\leq| & $\leq$ \\
\verb|\ll| & $\ll$ &
\verb|\mid| or \verb"|" & $\mid$ &
\verb|\models| & $\models$ \\ 
\verb|\neq| & $\neq$ &
\verb|\ni| & $\ni$ &
\verb|\notin| & $\notin$
\end{tabular}
\end{frame}
\begin{frame}[fragile]\frametitle{Symbols: relations\dots}
\begin{tabular}{llllll}
\verb|\parallel| & $\parallel$ &
\verb|\prec|& $\prec$ &
\verb|\preceq|& $\preceq$ \\
\verb|\perp| & $\perp$ &
\verb|\propto| & $\propto$ &
\verb|\sim| & $\sim$ \\
\verb|\simeq| & $\simeq$ &
\verb|\smile| & $\smile$ &
\verb|\sqsubseteq| & $\sqsubseteq$ \\
\verb|\sqsupseteq| & $\sqsupseteq$ &
\verb|\subset| & $\subset$ &
\verb|\subseteq| & $\subseteq$ \\
\verb|\succ| & $\succ$ &
\verb|\succeq| & $\succeq$ &
\verb|\supset| & $\supset$ \\
\verb|\supseteq| & $\supseteq$ &
\verb|\vdash| & $\vdash$ &
\end{tabular}
\end{frame}


\begin{frame}[fragile]\frametitle{Symbols: binary operators}
\begin{tabular}{llll}
\verb|\amalg| & $\amalg$ &
\verb|\ast| & $\ast$ \\
\verb|\bullet| & $\bullet$ &
\verb|\bigcirc| & $\bigcirc$ \\
\verb|\bigtriangledown| & $\bigtriangledown$ &
\verb|\bigtriangleup| & $\bigtriangleup$ \\
\verb|\cap| & $\cap$ &
\verb|\cdot| & $\cdot$ \\
\verb|\circ| & $\circ$ &
\verb|\cup| & $\cup$  \\
\verb|\dagger| & $\dagger$ &
\verb|\ddagger| & $\ddagger$ \\
\verb|\diamond| & $\diamond$ &
\verb|\div| & $\div$ \\
\verb|\mp| & $\mp$ &
\verb|\odot| & $\odot$ 
\end{tabular}
\end{frame}
\begin{frame}[fragile]\frametitle{Symbols: binary operators\dots}
\begin{tabular}{llll}
\verb|\ominus| & $\ominus$ &
\verb|\oplus| & $\oplus$ \\
\verb|\oslash| & $\oslash$ &
\verb|\otimes| & $\otimes$ \\
\verb|\pm| & $\pm$ &
\verb|\setminus| & $\setminus$ \\
\verb|\sqcap| & $\sqcap$ &
\verb|\sqcup| & $\sqcup$ \\
\verb|\star| & $\star$ &
\verb|\times| & $\times$ \\
\verb|\triangleleft| & $\triangleleft$ &
\verb|\triangleright| & $\triangleright$ \\
\verb|\uplus| & $\uplus$ &
\verb|\vee| & $\vee$ \\
\verb|\wedge| & $\wedge$ &
\verb|\wr| & $\wr$
\end{tabular}
\end{frame}


\begin{frame}[fragile]\frametitle{Symbols: arrows}
\begin{tabular}{llll}
\verb|\downarrow| & $\downarrow$ \\
\verb|\Downarrow| & $\Downarrow$ \\
\verb|\hookleftarrow| & $\hookleftarrow$ \\
\verb|\hookrightarrow| & $\hookrightarrow$ \\
\verb|\leftarrow| or \verb|\gets| & $\leftarrow$ \\
\verb|\Leftarrow| & $\Leftarrow$ \\
\verb|\leftharpoondown| & $\leftharpoondown$ \\
\verb|\leftharpoonup| & $\leftharpoonup$ \\
\verb|\leftrightarrow| & $\leftrightarrow$ \\
\verb|\Leftrightarrow| & $\Leftrightarrow$ 
\end{tabular}
\end{frame}
\begin{frame}[fragile]\frametitle{Symbols: arrows\dots}
\begin{tabular}{llll}
\verb|\longleftarrow| & $\longleftarrow$ \\
\verb|\Longleftarrow| & $\Longleftarrow$ \\
\verb|\longleftrightarrow| & $\longleftrightarrow$ \\
\verb|\Longleftrightarrow| & $\Longleftrightarrow$ \\
\verb|\longmapsto| & $\longmapsto$ \\
\verb|\longrightarrow| & $\longrightarrow$ \\
\verb|\Longrightarrow| & $\Longrightarrow$ \\
\verb|\mapsto| & $\mapsto$ \\
\verb|\nearrow| & $\nearrow$ \\
\verb|\nwarrow| & $\nwarrow$ \\
\end{tabular}
\end{frame}
\begin{frame}[fragile]\frametitle{Symbols: arrows\dots}
\begin{tabular}{llll}
\verb|\rightarrow| or \verb|\to| & $\to$ \\
\verb|\Rightarrow| & $\Rightarrow$ \\
\verb|\rightharpoondown| & $\rightharpoondown$ \\
\verb|\rightharpoonup| & $\rightharpoonup$ \\
\verb|\rightleftharpoons| & $\rightleftharpoons$ \\
\verb|\searrow| & $\searrow$ \\
\verb|\swarrow| & $\swarrow$ \\
\verb|\uparrow| & $\uparrow$ \\
\verb|\Uparrow| & $\Uparrow$ \\
\verb|\updownarrow| & $\updownarrow$ \\
\verb|\Updownarrow| & $\Updownarrow$ \\
\end{tabular}
\end{frame}


\begin{frame}[fragile]\frametitle{Symbols: big symbols}
\begin{tabular}{llllll}
\verb|\sum| & $\sum$ &
\verb|\int| & $\int$ &
\verb|\oint| & $\oint$ \\[5pt]
\verb|\prod| & $\prod$ &
\verb|\coprod| & $\coprod$ &
\verb|\bigcap| & $\bigcap$ \\[5pt]
\verb|\bigcup| & $\bigcup$ &
\verb|\bigsqcup| & $\bigsqcup$ &
\verb|\bigvee| & $\bigvee$ \\[5pt]
\verb|\bigwedge| & $\bigwedge$ &
\verb|\bigodot| & $\bigodot$ &
\verb|\bigotimes| & $\bigotimes$ \\[5pt]
\verb|\bigoplus| & $\bigoplus$ &
\verb|\biguplus| & $\biguplus$
\end{tabular}
\end{frame}


\end{document}
